\section{Claims, stops, and the next arithmetic door}
\label{sec:claim-boundary}

\subsection{L0 results proved here}

For finite atomic-normalized native incidences:
\begin{enumerate}[label=\textup{(\arabic*)},leftmargin=3em]
 \item the same-residual and erasure defects decompose exactly over complete
       native keys:
       \[
       R=\sum_iR_i,\qquad E=\sum_iE_i;
       \]
 \item \(R_i\) is residual-block diagonal and \(E_i\) is off block, giving
       \[
       \tr(E_iP(R_i))=0;
       \]
 \item the first mixed trace is the cross-key rectangle sum
       \[
       \tr(ER)=\sum_{i\ne j}\tr(E_iR_j);
       \]
 \item the rectangle coefficient is invariant under independent row and
       native-key unit gauges;
 \item every erasure-count sector \(S_{k,b}\) has the exact cyclic
       composition formula \eqref{eq:cyclic-sector-formula};
 \item the complete two-erasure sector in every even moment is nonnegative,
       with the divided-difference formula
       \eqref{eq:quadratic-divided-difference}.
\end{enumerate}

\subsection{L1 identifications}

The matrices above are identified with the actual TPC fixed-shift
interface by retaining:
\begin{itemize}
 \item the one fixed nonzero shift \(h_0\);
 \item every target \(mj+h_0\), largest prime \(p\), cofactor \(r\), and
       residual label \(s=d_mr\);
 \item every complete native key, including side, external block,
       opposite row, \(\gamma\), and physical orbit time;
 \item every literal complex coefficient \(\beta_w\);
 \item the raw row mass \(W_M(m)\); and
 \item every actual missing and first-failure condition.
\end{itemize}
On that physical squarefree carrier, the exact source-row gauge gives two
additional L1 identities: source-row signs telescope on every decorated
closed walk, leaving the odd residual boundary
\eqref{eq:odd-residual-boundary}, and a one-erasure word retains exactly
one unequal-residual sign pair.  The rectangle reduction
\eqref{eq:rectangle-sign-reduction} is their shortest instance.
The direct-cell \(\Gamma\)-formula is L1 only under its explicitly declared
identical-support specialization.  The general atom formula is the
unconditional L1 object.

\subsection{L2 statements not proved}

This paper does not prove:
\begin{itemize}
 \item cancellation in one literal cross-key rectangle sum;
 \item a uniform local compensation estimate in every active row;
 \item a saving for
       \(2q\tr(ER^{2q-1})\) at fixed or growing \(q\);
 \item an upper bound of the required scale for the nonnegative sector
       \(S_{2q,2}\);
 \item a complete even-moment bound uniform in all physical blocks;
 \item \(\|R+E\|=X^{o(1)}\), much less \(\|R+E\|=o(1)\);
 \item the missing-mass and represented-frame estimates in
       \eqref{eq:missing-mass-exponent}--%
       \eqref{eq:represented-lower-exponent};
 \item a negligible shorting penalty \(C_N\); or
 \item any parity-breaking estimate, prime-pair theorem, or twin-prime
       conclusion.
\end{itemize}

\subsection{Kill criteria}

\begin{enumerate}[label=\textup{(K\arabic*)},leftmargin=3.5em]
 \item \textbf{Incomplete key.}  A calculation that removes side,
       external block, opposite row, \(\gamma\), time, or another
       still-orthogonal output label estimates a different operator.

 \item \textbf{False one-key door.}  Searching for
       \(\tr(E_iR_i)\) cancellation cannot advance the mixed second moment:
       that trace is identically zero.  The first target is cross-key.

 \item \textbf{Wrong parity statistic.}  The number of erasure edges does
       not determine the sign.  Only the odd residual boundary does.

 \item \textbf{Quadratic negative-correction stop.}  A proposal that uses
       the complete \(S_{2q,2}\) as a negative correction is impossible:
       \(S_{2q,2}\ge0\).  This does not kill the linear or higher sectors.

 \item \textbf{Global-to-local stop.}  A small \(\tr(ER)\), an averaged
       row estimate, or a normalized trace does not prove the uniform local
       star bound.

 \item \textbf{Arithmetic scope mismatch.}  A fixed-pair theorem, an
       average over keys or shifts, a density-one statement, or a random
       sign model does not give the required complete growing-family
       estimate at the prescribed fixed \(h_0\).

 \item \textbf{Moment nonuniformity.}  If \(q=q_X\) grows, estimates whose
       constants depend uncontrolledly on \(q\) do not imply
       \eqref{eq:moment-operator-bound}.

 \item \textbf{Reassembly mismatch.}  The missing-native defect \(R+E\) is
       not \(H_{\rm dir}\), and neither controls \(C_N\) without additional
       physical comparison.

 \item \textbf{Endpoint stop.}  If the combined literal loss is at least
       \(1/400\), the present endpoint route must stop even if every
       intermediate theorem is correct.
\end{enumerate}

\subsection{The next specified door}

The shortest unresolved arithmetic object is now
\[
 \tr(ER)
\]
in the exact rectangle form
\eqref{eq:direct-cell-two-key-rectangle}.  A useful next step should:
\begin{enumerate}[label=\textup{(\roman*)},leftmargin=3em]
 \item census the actual two-key rectangles without compressing their
       native labels;
 \item split them by the two determinant equations and residual
       equality/inequality pattern;
 \item determine whether the rectangle holonomy admits a literal pairing,
       dispersion identity, or local star bound; and
 \item either prove a uniform saving or construct an actual-carrier
       obstruction closing this linear door.
\end{enumerate}
If the linear sector is structurally too large or its residual boundary
freezes on the actual support, the next alternatives are higher
one-erasure polygons, literal complex-phase cancellation, or direct
shorted reassembly.  The nonnegative quadratic sector is not an alternative
source of negative compensation.

\subsection{Status line}

TPC--69 moves the signed operator route from edgewise Möbius products to the
first exact cross-key holonomy and proves a new positivity obstruction for
the quadratic erasure response.  This is a genuine L0/L1 refinement of the
proof architecture.  It is not an L2 cancellation theorem and provides no
evidence that the parity barrier has been broken.
